\documentclass{beamer}

\usepackage[utf8]{inputenc}
\usepackage[brazil]{babel}
\usepackage{amsmath, amssymb}
\usepackage{physics}

\title{Computação Quântica de Variáveis Contínuas aplicada à Otimização Combinatória}
\author{Clovis Aparecido Caface Filho}
\institute{UFABC}
\date{2026}

\begin{document}

\frame{\titlepage}

%------------------------------------------------

\begin{frame}{Introdução}
\begin{itemize}
    \item Computação clássica baseada no modelo de Turing
    \item Problemas combinatórios possuem alta complexidade
    \item Aplicações:
    \begin{itemize}
        \item Logística
        \item Transporte
        \item Planejamento de rotas
    \end{itemize}
    \item Problemas como VRP são NP-difíceis
\end{itemize}
\end{frame}

%------------------------------------------------

\begin{frame}{Problema de Roteamento de Veículos (VRP)}
\begin{itemize}
    \item Generalização do Problema do Caixeiro Viajante (TSP)
    \item Múltiplos veículos e restrições:
    \begin{itemize}
        \item Capacidade
        \item Tempo
        \item Custos
    \end{itemize}
    \item Amplamente utilizado em logística e transporte
\end{itemize}
\end{frame}

%------------------------------------------------

\begin{frame}{Formulação Matemática do VRP}
\begin{itemize}
    \item Grafo completo $G=(V,E)$
    \item Função objetivo:
\end{itemize}

$$
\min \sum_{i \in V} \sum_{j \in V, j \neq i} c_{ij} x_{ij}
$$

\begin{itemize}
    \item $x_{ij}$ indica se a aresta é utilizada
    \item Restrições garantem:
    \begin{itemize}
        \item Visita única
        \item Conservação de fluxo
        \item Capacidade
    \end{itemize}
\end{itemize}
\end{frame}

%------------------------------------------------

\begin{frame}{Computação Quântica}
\begin{itemize}
    \item Informação representada por qubits
    \item Superposição:
\end{itemize}

$$
\ket{\psi} = \alpha \ket{0} + \beta \ket{1}
$$

\begin{itemize}
    \item Propriedades:
    \begin{itemize}
        \item Superposição
        \item Emaranhamento
    \end{itemize}
    \item Possível vantagem sobre métodos clássicos
\end{itemize}
\end{frame}

%------------------------------------------------

\begin{frame}{Variáveis Contínuas (CVQC)}
\begin{itemize}
    \item Informação codificada em variáveis contínuas
    \item Baseada em sistemas bosônicos (fótons)
    \item Estados definidos por:
\end{itemize}

\[
\ket{\psi} = \int \psi(x)\ket{x} \text{d}x
\]

\begin{itemize}
    \item Espaço de Hilbert infinito
\end{itemize}
\end{frame}

%------------------------------------------------

\begin{frame}{Operadores em CVQC}
\begin{itemize}
    \item Deslocamento:
\end{itemize}

\[
D(\alpha) = \exp(\alpha \hat{a}^\dagger - \alpha^* \hat{a})
\]

\begin{itemize}
    \item Compressão (squeezing)
    \item Rotação
    \item Portas não gaussianas (Kerr, cúbica)
\end{itemize}
\end{frame}

%------------------------------------------------

\begin{frame}{Objetivos}
\begin{itemize}
    \item Desenvolver modelo quântico para otimização combinatória
    \item Aplicação em VRP
    \item Objetivos específicos:
    \begin{itemize}
        \item Estudo do TSP
        \item Formulação matemática compatível com CVQC
        \item Mapeamento em Hamiltonianos
    \end{itemize}
\end{itemize}
\end{frame}

%------------------------------------------------

\begin{frame}{Conclusão}
\begin{itemize}
    \item VRP é um problema relevante e complexo
    \item Computação quântica oferece novas abordagens
    \item CVQC é promissora para problemas contínuos
    \item Integração entre matemática e física é essencial
\end{itemize}
\end{frame}

%------------------------------------------------

\end{document}